\section{Consequences for the rotation experiment}\label{sec:rotation34}
We now apply the budget to the original wordwise specification, not to a
statewise-intertwining subclass.

\begin{proposition}[Terminal calibration of the past phase]\label{prop:terminal34}
A realization of \eqref{eq:array34} with uniform row-TV error at most
$\varepsilon<1/20$ satisfies $q_N\ge1/10-2\varepsilon$ under the fair-command
test law with fixed seed $x=(1,1)$.
\end{proposition}
\begin{proof}
Let $d=d_1+i d_2$ be the complex conditional mean of the two possible
answers from the final state. Its modulus is at most $\sqrt2$. For any
fixed command word, exactness would give
\[
 \E[d(S_N)\mid C_1,\ldots,C_N]
                  =\frac{1+i}{10}\zeta^{L_N}.
\]
A TV error $\varepsilon$ for a binary answer changes its mean by at most
$2\varepsilon$. The complex error in this display is therefore at most
$2\sqrt2\varepsilon$. Multiply by $Y_N=\zeta^{-L_N}$ and average. The reverse
triangle inequality gives
\[
 |\E[Y_Nd(S_N)]|\ge\sqrt2(1/10-2\varepsilon).
\]
Apply Lemma~\ref{lem:phase34}. Only the final-length answer identities
were used. There is no assumption of a correct answer at an earlier cut.
\end{proof}

\begin{proof}[Proof of the lower assertions of Theorem~\ref{thm:main34}]
Proposition~\ref{prop:terminal34} and Theorem~\ref{thm:fourier34} give
\eqref{eq:main-count34}. If every cut has at most $K$ states, then
$N\beta_{2K}(\alpha)\le18K^3/\kappa^2$. For $b_\alpha>0$,
\[
 |\sin(\pi l\alpha)|\ge2\|l\alpha\|_{\R/\mathbb Z}
          \ge\frac{2b_\alpha}{l}\ge\frac{b_\alpha}{K}
          \quad(1\le l\le2K).
\]
This proves \eqref{eq:main-lower34}.

For the golden angle set $\alpha'=-(1+\sqrt5)/2$. If $p$ is an integer
nearest $l\alpha$, then
\[
 (l\alpha-p)(l\alpha'-p)=p^2+lp-l^2\in\mathbb Z\setminus\{0\},
 \qquad |l\alpha'-p|\le\sqrt5l+1/2<3l.
\]
Hence $l\|l\alpha\|>1/3$ and $b_\alpha\ge1/3$. With $\kappa=1/10$
this gives $(N/16200)^{1/5}$. The exact rank-three bound is
Proposition~\ref{prop:separate34}.

For $\zeta=(3+4i)/5$, no positive power equals one: otherwise
$\zeta+\zeta^{-1}=6/5$ would be a rational algebraic integer. Thus
$(3+4i)^l-5^l$ is a nonzero Gaussian integer and has modulus at least one.
It follows that
\[
 \sin^2(\pi l\alpha)=\frac{|\zeta^l-1|^2}{4}\ge\frac{25^{-l}}4,
 \qquad \beta_{2K}(\alpha)\ge\frac{25^{-2K}}4.
\]
The profile bound gives $N\le7200K^3 25^{2K}$ at zero error. Since
$K^3\le8^K$ for integer $K\ge1$, it also gives
$N\le7200\,5000^K$. This proves \eqref{eq:rational-head34}.
\end{proof}

The occupation conclusion contains more information than a peak bound.
For example, at the golden angle and zero error there are at most
$16200k^5$ command epochs with at most $k$ available labels, irrespective
of the widths at the other epochs. The analogous rational-angle count is
at most $7200k^3 25^{2k}$. These are proof bounds, not exact finite
complexity tables.

\subsection{An exact upper construction retained from the preceding argument}
For completeness we give the construction needed for the upper side of
\eqref{eq:golden34}; its role here is inherited \cite{GTF33}. A convergent
$p/q$ of $\alpha$ satisfies $|\alpha-p/q|<q^{-2}$. Choose one with
$q\ge\max\{5,(40N)^{1/3}\}$ and put
\[
 a=\pi/q,\qquad \delta=2\pi(\alpha-p/q),\qquad
 \lambda=\frac{\cos(a-|\delta|)}{\cos a},\qquad
 \omega=\frac{\sin|\delta|}{\lambda\sin(2a)}.
\]
Here $|\delta|<2\pi/q^2<a$, so $\lambda\ge1$ and $0\le\omega\le1$.
Write $u_s=(\cos(2\pi s/q),\sin(2\pi s/q))$ for the regular polygon
vertices, with subscripts modulo $q$.

\begin{proposition}[Resonant stochastic polygon]\label{prop:upper34}
There is a $q$-state exact realization of \eqref{eq:array34}, with command
matrices independent of the epoch and a horizon-dependent final decoder.
For bounded partial quotients it gives $W_N=O_\alpha(N^{1/3})$.
For the golden angle the upper bound in \eqref{eq:golden34} holds.
\end{proposition}
\begin{proof}
At cut $t$ use the vector $r_tu_s$ for label $s$, where
$r_t=\lambda^t/4$. The idle row is
\[
 T_0(s,\cdot)=\lambda^{-1}\delta_s+(1-\lambda^{-1})\operatorname{Unif}[q].
\]
The active row is
\[
 T_1(s,\cdot)=(1-\omega)\delta_{s+p}
                  +\omega\delta_{s+p+\operatorname{sgn}\delta}.
\]
If $\delta=0$ use $\omega=0$ and omit the second term. The line joining
$u_{s+p}$ to its indicated neighbor intersects the ray at angle
$2\pi s/q+2\pi\alpha$ at radius $1/\lambda$. Resolving its two
coordinates, or using the sine rule, gives exactly the displayed $\omega$.
Consequently the expected next vector is $r_tu_s$ after an idle command
and $R_\alpha r_tu_s$ after an active command.

Since $\tan(\pi/q)\le2\pi/q$ for $q\ge5$,
\[
 \log\lambda=\int_{a-|\delta|}^a\tan v\,dv
        \le|\delta|\tan a<4\pi^2/q^3<40/q^3.
\]
Thus $r_N\le e/4<1$. The polygon at $t=0$ has inradius
$\cos(\pi/q)/4>\sqrt2/10$ and therefore contains all four seed vectors
$x/10$. Initialize by a convex representation of the seed in that polygon.
The final decoder uses mean $r_Ne_j^Tu_s$, which lies in $[-1,1]$.
Induction on the commands proves exactness for every word.

Consecutive convergent denominators have bounded ratios at bounded-type
angles, yielding the claimed order. At the golden angle the denominators
are Fibonacci numbers, and the first one not smaller than a threshold
$L\ge5$ is strictly less than $2L$. This gives the stated numerical bound.
\end{proof}

For the rational rotation an elementary exact rational alternative retains
a triangle label and the number of active commands. Initialize in the
triangle from Proposition~\ref{prop:separate34}, increment the count only
on the active command, and decode the rotated triangle vertex. It uses at
most $3(N+1)$ labels and rational rows. The preceding resonant construction
uses exact real rows; a finite fair-bit implementation of those rows is
not being asserted.
